% SGA 3，第 VII_B 篇，英文重建稿——第 1 轮。
% 范围：权威底本分篇第 92 页，第 5.2.2 节全文，止于第 5.2.3 节之前。
% 合并阅读本页码： 605.
% 再版印刷页码： 594.
% 本范围内原版页边定位： 559.
% 权威底本： Polo--Gille Exp7B-13oct24.pdf,
% SHA-256 D303690E0195654BBE7C3DAA5D89EA23E24CD709A3575217234AC19AE8836694.
% 曾参考 Jacob Reinhold 的 CC BY 4.0 Markdown 作为比对谱系。
% 本组件没有图表；全部数学均为原生 LaTeX。

\subsubsection*{5.2.2}
\phantomsection
\label{sga3:VIIB:5.2.2}

先假设 \(B\) 的高度至多为 \(1\)；即
\[
  x^p=0
  \qquad\text{对每个 }x\in I.
\]
由 \SGARefLink{sga3:VIIB:5.1.4}{5.1.4}，与 \(B\) 的滤过
\[
  I\supset I^2\supset I^3\supset\cdots
\]
相联系的分次代数
\[
  \operatorname{Gr}_I(B)
\]
带有范畴 \(\mathbf{Alpg}/k\) 中的余群结构。换言之，预有限代数
\[
  C=\prod_{n\geq0}\operatorname{Gr}_I(B)_n
\]
是某个形式 \(k\)-群 \(N\) 的仿射代数。显然，
\[
  \omega_{N/k}=I/I^2,
\]
且 \(N\) 是高度至多为 \(1\) 的无穷小群。

因此，由 \SGARefLink{sga3:viib:subsection:4-4}{4.4}，\(\omega_{N/k}\) 的恒等映射诱导同构
\[
  B'
  =
  k[[\omega_{N/k}]]
  \big/
  \bigl((x^p)\bigr)_{x\in\omega_{N/k}}
  \xrightarrow{\ \sim\ }
  C.
\]
特别地，当 \(B'\) 和 \(B\) 都按其增广理想的幂赋予滤过时，\SGARefLink{sga3:VIIB:5.2.1:section-proposition}{5.2.1} 的映射 \(\phi\) 在相联系的分次代数之间诱导同构。由 \textup{[CA]}，\(\S\)V.7，引理~1\textup{（亦见 \SGARefLink{sga3:VIIB:5.1.5}{5.1.5}）}，可知 \(\phi\) 为同构。
